\documentclass[12pt,a4paper]{article}
\usepackage{preamble}

\title{\textbf{Solusi dan Kunci Jawaban Ujian Akhir Semester Matematika Diskrit}}
\author{Kevin Wirya Valerian - 13524019}
\date{}

\begin{document}

\SetWatermarkText{}
\SetWatermarkScale{3}

\maketitle

\section*{Bagian I: Easy}

\subsection*{Soal 1 [Logika Penalaran -- Inferensi Logika]}
\textbf{Solusi:} \\
Misalkan proposisi yang ada dalam bentuk variabel.
\begin{itemize}
    \item $p$: Koro-sensei pergi ke bulan.
    \item $q$: Koro-sensei membawa oleh-oleh batu bulan.
    \item $r$: Koro-sensei menonton Netflix di rumah.
\end{itemize}

Pernyataan-pernyataan diubah menjadi bentuk proposisi majemuk.
\begin{enumerate}
    \item $p \rightarrow q$ (Premis 1)
    \item $\neg q \lor r \equiv q \rightarrow r$ (Premis 2)
    \item $p \land \neg r$ (Premis 3)
\end{enumerate}

\textbf{Langkah Inferensi:}
\begin{itemize}
    \item Dari premis (3) $p \land \neg r$, menggunakan hukum simplifikasi diperoleh $p$ dan $\neg r$.
    \item Dari $p$ dan premis (1) $p \rightarrow q$, menggunakan Modus Ponens, diperoleh $q$.
    \item Dari $q$ dan premis (2) $q \rightarrow r$, menggunakan Modus Ponens, diperoleh $r$.
    \item Terjadi kontradiksi antara kesimpulan $r$ dengan fakta dari premis (3) bahwa $\neg r$.
\end{itemize}
$\therefore$ Rencana Koro-sensei \fbox{tidak konsisten} karena menghasilkan kontradiksi.

\subsection*{Soal 2 [Himpunan -- Kardinalitas]}
\textbf{Solusi:} \\
Misalkan himpunan siswa yang menyukai cabang olahraga:
\begin{itemize}
    \item $S$: Sepak bola dengan $|S| = 35$
    \item $T$: Bulu tangkis dengan $|T| = 50$
    \item $B$: Basket dengan $|B| = 40$
\end{itemize}
Diketahui pula data irisan dan siswa yang tidak menyukai ketiganya:
\[ |S \cap T| = 21, \quad |T \cap B| = 17, \quad |S \cap T \cap B| = 9 \]
Total siswa $= 100$, dan siswa yang tidak menyukai ketiganya $= 8$. Maka, banyak siswa yang menyukai minimal satu olahraga adalah:
\[ |S \cup T \cup B| = 100 - 8 = 92 \]

\begin{enumerate}[label=(\alph*)]
    \item Banyak siswa yang senang bermain sepak bola dan basket dinyatakan dalam notasi himpunan sebagai
    \[ |S \cap B| \]
    
    \item Berdasarkan prinsip inklusi-eksklusi untuk 3 himpunan,
    \[ |S \cup T \cup B| = |S| + |T| + |B| - |S \cap T| - |T \cap B| - |S \cap B| + |S \cap T \cap B| \]
    Substitusikan nilai yang diketahui ke dalam persamaan.
    \begin{align*}
        92 &= 35 + 50 + 40 - 21 - 17 - |S \cap B| + 9 \\
        92 &= 96 - |S \cap B| \\
        |S \cap B| &= 96 - 92 = 4
    \end{align*}
\end{enumerate}
$\therefore$ Banyak siswa yang senang bermain sepak bola dan basket adalah \fbox{4 siswa}.

\subsection*{Soal 3 [Relasi dan Fungsi -- Sifat Relasi]}
\textbf{Solusi:} \\
Himpunan menu makanan $A = \{N, M, C, B\}$. Relasi $R$ didefinisikan dengan $(x, y) \in R$ apabila $x$ lebih disukai daripada $y$.

\begin{enumerate}[label=(\alph*)]
    \item Berdasarkan data preferensi:
    \begin{itemize}
        \item $N$ lebih disukai daripada $C$ dan $B \implies (N, C), (N, B) \in R$
        \item $B$ lebih disukai daripada $M$ dan $C \implies (B, M), (B, C) \in R$
        \item $M$ lebih disukai daripada $N \implies (M, N) \in R$
    \end{itemize}
    Maka relasi $R$ dalam bentuk himpunan pasangan terurut adalah
    \[ R = \{(N, C), (N, B), (B, M), (B, C), (M, N)\} \]

    \item Analisis sifat relasi $R$
    \begin{itemize}
        \item Refleksif: \textbf{Tidak}. Karena $(x, x) \notin R$ untuk setiap $x \in A$ (misalnya $(N, N) \notin R$).
        \item Setangkup: \textbf{Tidak}. Karena terdapat $(N, C) \in R$ tetapi $(C, N) \notin R$.
        \item Menghantar: \textbf{Tidak}. Karena terdapat $(N, B) \in R$ dan $(B, M) \in R$, tetapi $(N, M) \notin R$.
    \end{itemize}
\end{enumerate}

\vspace{0.5cm}

\subsection*{Soal 4 [Aljabar Boolean -- Bentuk Kanonik SOP dan POS]}
\textbf{Solusi:} \\
Berdasarkan tabel kebenaran fungsi Boolean $f(x, y, z)$:

\begin{enumerate}[label=(\alph*)]
    \item Fungsi bernilai 1 pada baris minterm $m_0, m_2, m_5, m_6$:
    \begin{itemize}
        \item $m_0 (0,0,0) \rightarrow x'y'z'$
        \item $m_2 (0,1,0) \rightarrow x'yz'$
        \item $m_5 (1,0,1) \rightarrow xy'z$
        \item $m_6 (1,1,0) \rightarrow xyz'$
    \end{itemize}
    Maka bentuk kanonik SOP adalah
    \[ f(x, y, z) = \sum(0, 2, 5, 6) = x'y'z' + x'yz' + xy'z + xyz' \]

    \item Fungsi bernilai 0 pada baris maxterm $M_1, M_3, M_4, M_7$:
    \begin{itemize}
        \item $M_1 (0,0,1) \rightarrow (x + y + z')$
        \item $M_3 (0,1,1) \rightarrow (x + y' + z')$
        \item $M_4 (1,0,0) \rightarrow (x' + y + z)$
        \item $M_7 (1,1,1) \rightarrow (x' + y' + z')$
    \end{itemize}
    Maka bentuk kanonik POS adalah
    \[ f(x, y, z) = \prod(1, 3, 4, 7) = (x + y + z')(x + y' + z')(x' + y + z)(x' + y' + z') \]
\end{enumerate}

\vspace{0.5cm}

\subsection*{Soal 5 [Kompleksitas Algoritma -- Notasi Big-O]}
\textbf{Solusi:}

\begin{enumerate}[label=(\alph*)]
    \item Berikut adalah analisis kompleksitas waktu untuk masing-masing algoritma.
    \begin{itemize}
        \item Algoritma A
        \[ T_A(n) = (n + 4)^2 - 8n = n^2 + 8n + 16 - 8n = n^2 + 16 \implies O(n^2) \]
        
        \item Algoritma B
        \[ T_B(n) = n \log_2(n^5) + n^{1.5} = 5n \log_2 n + n^{1.5} \]
        Karena $n^{1.5} = n\sqrt{n}$ tumbuh lebih cepat daripada $n \log_2 n$ untuk $n$ besar, $T_B(n) = O(n^{1.5})$.
        
        \item Algoritma C
        \[ T_C(n) = 4(\log_2 n)^3 + 3n \]
        Suku yang lebih dominan adalah $3n$ sehingga $T_C(n) = O(n)$.
        
        \item Algoritma D
        \[ T_D(n) = 2^{\log_2(n \log_2 n)} + \log_2 n = n \log_2 n + \log_2 n \]
        Suku dominan adalah $n \log_2 n$ sehingga $T_D(n) = O(n \log n)$.
        
        \item Algoritma E
        \[ T_E(n) = 10(\log_2 n)^2 + 50 \implies O((\log n)^2) \]
    \end{itemize}

    \item Urutan laju pertumbuhan fungsi dari yang paling efisien hingga paling tidak efisien adalah
    \[ O((\log n)^2) < O(n) < O(n \log n) < O(n^{1.5}) < O(n^2) \]
    Sehingga urutan algoritmanya adalah \fbox{$E<C<D<B<A$}
\end{enumerate}

\end{document}
